% Section 5 — Discussion
\section{Discussion}
\label{sec:discussion}

\subsection{Forecast accuracy is not decision quality}
\label{sec:d-acc-dec}
Accuracy metrics are symmetric, per-unit, and policy-free; inventory costs
are asymmetric, thresholded, and policy-defined. Safety stock converts error
volatility into money; the order-up-to rule converts error bias into money;
and the $P/H$ ratio sets the exchange rate. An evaluation that stops at the
forecast therefore measures only part of a model's economic effect.

\subsection{Why LSTM wins forecasting but not always inventory}
\label{sec:d-lstm}
The global LSTM captures patterns unavailable to per-series linear models:
zero-inflated regimes, cross-item weekly structure, and nonlinear rebounds.
That accounts for its forecast lead on both datasets. Inventory, however,
purchases neither point forecasts nor patterns directly. It purchases a
dispersion summary (error volatility, via safety stock) and a level (the
next-$L$ sum). On dense Store demand, the LSTM's sharpness includes
occasional confident under-predictions: inexpensive under MAE, expensive at
$P=5$. The moving average, by construction, rarely under-shoots by large
amounts. The relevant distinction is therefore between precision and caution
under asymmetric cost, not between intelligence and simplicity.

\subsection{Why simple models remain competitive}
\label{sec:d-simple}
Moving Average and SES satisfy two conditions the policy rewards. First,
averaging is close to optimal when most daily variation is unpredictable
noise; the recent mean is then a minimum-variance guess. Second, their
errors are approximately symmetric and stable, which is what safety-stock
formulas assume. Simplicity aligns with the policy's assumptions.
Sophistication improves on it only where the demand signal is strong enough
to repay the departure, which holds on sparse M5 demand but not on dense
Store demand under the default cost ratio.

\subsection{Demand structure and model suitability}
\label{sec:d-structure}
Intermittency determines which machinery is legitimate. Where zeros
dominate, methods that decline to fit them (the Croston family, the
Naive-fallback ARIMA guardrail) and learners that pool across series (the
LSTM) perform best, while linear autoregression lacks estimable structure.
Where demand is dense and weekly, seasonal specifications (SARIMA, Seasonal
Naive) contribute and trend specifications (DES, TES) overfit noise into
stock. The archetype tables constitute the suitability map; the Store-M5
contrast corroborates it.

\subsection{Inventory policy as the effective loss function}
\label{sec:d-policy-loss}
The policy parameters $(L, z, P/H)$ function as the effective loss: forecasters
are implicitly optimized for symmetric error but scored on asymmetric cost.
The sensitivity grid makes this visible, with the lowest-cost model changing
as $P$ moves from 3 to 10. Model selection should therefore be conditional on
the policy, and accuracy-only leaderboards should be read with that
conditionality stated.

\subsection{Statistical significance vs.\ practical importance}
\label{sec:d-sig}
At $n=112{,}000$ paired errors, significance is inexpensive and consequently
uninformative on its own. The M5 LSTM-TSB pair ($p \approx 2\times 10^{-10}$,
$d_z=-0.031$) is significant in the formal sense and negligible for any
business purpose. Effect sizes, intervals, and the MA-SES test disagreement
are what prevent overinterpretation. This qualification is strengthened, not
weakened, by the dependence structure noted in section 2.9:
correlated observations inflate nominal precision further, which is an
additional reason to weight $d_z$ over $p$.

\subsection{Complexity vs.\ value}
\label{sec:d-complex}
The ladder's complexity order predicts neither the accuracy order nor the
cost order. The only complexity that pays is the neural kind, and only in
accuracy everywhere and in cost in one environment. The practical reading is
budgetary: a moving average costs nothing to operate and nearly leads; an
LSTM costs substantial engineering and leads conditionally.

\subsection{Robust vs.\ fragile conclusions}
\label{sec:d-robust}
Stability is a property of conclusions under policy variation, not of
models. M5-LSTM is stable (25 of 27, with losses confined to one explicable
region). Store-MA is policy-sensitive (9 of 27), as is every Store
alternative. Policy-sensitivity does not invalidate the default-policy
result: under the stated policy, Moving Average genuinely records the lowest
cost. It conditions the result: changing lead time or cost ratio requires
re-running the comparison before acting on it.

\subsection{What can and cannot be generalized}
\label{sec:d-general}
Transferable findings concern mechanisms and methods: the forecast-decision
gap and its cost-asymmetry mechanism, the significance-versus-size
discipline, the sensitivity-as-stability method, LSTM forecast performance
across both structures, and the TES/DES failure modes. Not transferable are
the specific board winners (policy- and structure-contingent), anything about
the Smooth archetype ($n=1$), anything about LLMs (unexecuted), and anything
beyond $h=28$ or univariate demand. The contribution is a map with stated
boundaries rather than a universal ranking.
